\slajdid{04-s013}
\begin{frame}[label=04-s013]
\gusttekst\setlength{\tabcolsep}{4pt}\renewcommand{\arraystretch}{1.02}
\begin{columns}[t,onlytextwidth]
\column{.47\textwidth}
PRIMER:\quad $13.2_{10}={?}_2$\par\vspace{2mm}
CELOBROJNA VREDNOST\par
\begin{tikzpicture}
\node[inner sep=0] (t) {\begin{tabular}{ccc}\toprule
CV:2 & REZULTAT & OSTATAK\\\midrule
13:2 &6&1\\6:2&3&0\\3:2&1&1\\
\textcolor{DEcrvena}{1:2}&\textcolor{DEcrvena}{0}&\textcolor{DEcrvena}{1}\\
0:2&0&0\\\bottomrule
\end{tabular}};
\draw[DEplava,->] ([xshift=3mm]t.south east)--([xshift=3mm,yshift=-5mm]t.north east);
\node[anchor=north east] at ([xshift=7mm,yshift=-1mm]t.south east) {1101};
\end{tikzpicture}
\column{.50\textwidth}
RAZLOMLJENA VREDNOST\par
\begin{tikzpicture}
\node[inner sep=0] (t) {\begin{tabular}{ccc}\toprule
RV x 2 & REZULTAT & CEO DEO\\\midrule
0.2 x 2&0.4&0\\0.4 x 2&0.8&0\\0.8 x 2&1.6&1\\0.6 x 2&1.2&1\\
0.2 x 2&0.4&0\\0.4 x 2&0.8&0\\0.8 x 2&1.6&1\\0.6 x 2&1.2&1\\0.2 x 2&0.4&0\\
\ldots&\ldots&\ldots\\\bottomrule
\end{tabular}};
\draw[DEplava,->] ([xshift=3mm]t.north east)--([xshift=3mm]t.south east);
\node[anchor=north east] at ([xshift=6mm,yshift=-1mm]t.south east) {.001100110\ldots.};
\end{tikzpicture}
\end{columns}
\[13.2_{10}=1101.001100110\ldots_2\]
\alert{Što veći broj cifara u razlomljenom delu imaćemo tačnije „rezultat“ prelaska.\\
Apsolutno tačan nikada. Moramo se dogovoriti koji broj cifara uzimamo.\\
U ovom primeru. Za neke druge primere mogli bi da dobijemo i tačan rezultat prelaska. Na primer 0.5 ili 0.25 ili 0.75 ili \ldots.}
\end{frame}
